\subsection{The exact secondary character of a global-unit norm}

\begin{statement}{A.184}{Proved}
{All-height relative-unit residues and the two-step norm image}
Use the fixed fields \(K_n,E_n\), the actual norms, and the unit
\(U=\alpha^2-2\) of A.181.  Put
\(\mathcal U_n=\mathcal O_{K_n}^{\times}\).
At each \(q=11,59\), the pure quadratic-residue prime has
residue field \(\F_{q^2}\) at every height.  Use its fixed
identification from \(K_1\) and define
\[
 \theta_n:\mathcal U_n\longrightarrow\F_3^2,\qquad
 \theta_{n,q}(u)=\log_{\lambda_q(\alpha)}\lambda_q(u),
 \quad\lambda_q(u)=\overline u^{(q^2-1)/3}.
\tag{A.184.1}
\]
The two generators \(\lambda_q(\alpha)\) have exact order three.
For every \(n\ge1\), the following maps and equalities hold:
\[
\begin{gathered}
 N_{n+1,n}\mathcal U_{n+1}=M_n:=\ker\theta_n,\\
 \theta_{n+1}\big|_{\ker(N_{n+1,n}:\mathcal U_{n+1}
                                    \to\mathcal U_n)}=0,\\
 \psi_n:M_n\longrightarrow\F_3^2,\qquad
       \psi_n(u)=\theta_{n+1}(v)\quad(N_{n+1,n}v=u),\\
 N_{n+2,n}\mathcal U_{n+2}=\ker(\psi_n:M_n\to\F_3^2),
 \qquad \psi_n(u^3)=\theta_n(u)\quad(u\in\mathcal U_n).
\end{gathered}
\tag{A.184.2}
\]
The third line is a well-defined homomorphism independent of the
chosen actual global-unit preimage.  All equalities, kernels and
the scalar identity in A.184.2 also hold after pro-three completion,
with the continuous extensions of \(\theta_n,\psi_n\) and the
actual completed norm maps.

In particular the scalar \(U^3\in K_n\) is a one-step actual
global-unit norm and is not a two-step completed global-unit norm,
at every \(n\).  It is an actual field norm at every finite
higher height and has compatible local norm families, while its
actual and completed global inverse field-norm fibres are empty.
The finite two-step global-unit support class has exact order three.
These statements retain the prescribed target \(U^3\) at the
indicated level.  A.184.2 gives the evaluated additional unit
equation for two consecutive steps; it assigns no stabilization
property to the images from arbitrarily distant levels.
\end{statement}

\begin{proof}
\emph{The one-step local and global norm maps.}
At the quadratic branch over \(q\), the local valuation of
\(\rho_n\) is eight.  The next field is a totally ramified
cubic, and its unit norms are exactly the units with cube residue.
Indeed residue norm is the third power, scalar lifts give every
cube residue, and cubing is an automorphism on the principal
units because \(q\ne3\).  Since
\(v_3(q^2-1)=1\), the kernel is exactly the corresponding
coordinate of A.184.1.  At the linear branch the residue field
has order \(q\equiv2\pmod3\), so every unit is a norm.
At three the pure field has odd degree over \(\Q_3\),
contains no \(\mu_3\), and has \(v(\rho_n)=1\).
The root-stabilizer inclusion in the abelianization of the
local holomorph is onto; the exact Artin/norm diagram in
A.140 gives surjectivity of the pure local field norm.
Total ramification makes a preimage of a unit a unit.
At every other finite place the root algebra is unramified,
and unit norms from its field factors are onto.  At a real
place choose its real factor, and at a complex place choose
one complex factor.  Thus the two coordinates in A.184.1
are precisely all local conditions on a global-unit target.

The global Hasse norm map is retained as well.  In
\(\Omega_{n+1}=\Q(\mu_{3^{n+1}},\rho_{n+1})\), the
groups fixing \(K_n\) and \(K_{n+1}\) are
\[
 G=(3^n\Z/3^{n+1})\rtimes(\Z/3^{n+1})^\times,
 \qquad H=(\Z/3^{n+1})^\times.
\tag{A.184.3}
\]
The decomposition subgroup over three is all of \(G\),
by the actual local/global degree calculation in A.162.
For the norm torus the lattice is
\(\Z[G/H]/\Z\sum gH\); restriction at three is the
identity in its decomposition-group \(H^2\) kernel.
The torus duality and inflation map of A.181.10 therefore
give the Hasse norm principle with precisely these groups.
Hence a global unit in \(\ker\theta_n\) is an actual
field norm.

The one-step global-unit repair proved in A.181 now removes
the support of that preimage.  Its cyclic step in
\(E_{n+1}/E_n\) uses a divisor \(D_q\) supported on
actual fixed changed-prime ideals, represents the three-primary
class of a split-orbit divisor, principalizes a prime-to-three
multiple, and uses Bezout to retain the exact target norm.
For the pure step, if the resulting unit in \(E_{n+1}\)
has norm \(u\), its quadratic norm \(v\) satisfies
\(N_{K_{n+1}/K_n}v=u^2\).  The actual unit \(v^2/u\)
then has norm \(u\), since its exponent is \(4-3=1\).
This proves the first equality in A.184.2, including signs.

\emph{Relative norm-one units and actual ramified divisor classes.}
Write \(L=E_{n+1}\), \(F=E_n\), let \(\sigma\)
generate \(L/F\) by \(\sigma\rho_{n+1}=\zeta_3\rho_{n+1}\),
and let \(\tau\) fix the positive pure fields and invert
\(\zeta_3\).  Put \(A_j=\operatorname{Cl}(E_j)\{3\}\).
The class computations A.173 and A.177 give, for the actual
ideal extension and norm,
\[
 \sigma|_{A_{n+1}}=1,\qquad
 N:A_{n+1}\twoheadrightarrow A_n,\qquad
 jA_n=3A_{n+1}.
\tag{A.184.4}
\]
For \(n=1\) this says \(A_1=0\) and \(3A_2=0\),
as proved in A.173.  For larger \(n\), ideal norm sends
each generating changed-prime class to its predecessor
with residue-degree weight one; the relation
\(jN=1+\sigma+\sigma^2=3\) proves the last equality.

Let \(u\in\mathcal U_{n+1}\) have pure relative norm
one.  Its norm in \(L/F\) is also one, and \(\tau u=u\).
Hilbert 90 gives an actual \(w\in L^\times\) with
\[
                         u=\sigma(w)/w.
\tag{A.184.5}
\]
Because \(u\) is a global unit, \(\operatorname{div}(w)\)
is \(\sigma\)-invariant.  At an unramified inert prime its
coefficient is the extension of the lower coefficient.  On
a split orbit its three coefficients are equal and likewise
are extended.  Only the seven ramified primes require residues
modulo their ramification index three.  Thus there is an
actual lower fractional divisor \(D\) with
\[
 \operatorname{div}(w)=jD+k_3R_3+
             \sum_{q\in\{11,59\}}\sum_{i=0}^2 k_{q,i}R_{q,i},
 \qquad k_3,k_{q,i}\in\{0,1,2\}.
\tag{A.184.6}
\]
Here \(\tau R_{q,0}=R_{q,0}\),
\(\tau R_{q,1}=R_{q,2}\), and \(\tau R_3=R_3\).
Reducing the class of this principal divisor in
\(A_{n+1}/3A_{n+1}\) removes \(jD\) by A.184.4.

At any changed-prime branch choose a lower uniformizer
\(\pi_0\).  Integers \(a,b\) with \(8a+3b=1\)
give the upper uniformizer
\(\pi=\rho_{n+1}^{a}\pi_0^b\), with
\(a\equiv2\pmod3\).  Its actual inertia action is
\[
                         \sigma(\pi)/\pi=\zeta_3^2.
\tag{A.184.7}
\]
The inertia action on a local unit's residue is trivial.
Writing \(w=\pi^k w_0\) therefore gives
\(\overline u=\zeta_3^{2k}\) at that branch.
At \(R_{q,0}\), the pure unit has residue in \(\F_q\),
which contains no nonidentity cube root of unity; hence
\(k_{q,0}=0\pmod3\).  At the exchanged two branches,
the identity \(\tau u=u\), together with
\(\tau\zeta_3=\zeta_3^{-1}\), gives
\[
                    k_{q,2}=-k_{q,1}\pmod3.
\tag{A.184.8}
\]

The two unramified cubic class characters of A.179 are the
base changes of the radicals
\(z_1=11(1+6\alpha)\) and \(z_2=649\).
Both radicals are \(\tau\)-fixed, so the Kummer pairing
with \(\mu_3\) shows that their ordinary class characters
\(\chi_1,\chi_2\) are \(\tau\)-anti-invariant.
They kill the \(\tau\)-fixed three-prime class and kill
\(3A_{n+1}=jA_n\).  On the differences
\(\Delta_q=[R_{q,1}]-[R_{q,2}]\), their literal matrix is
\[
 \begin{pmatrix}
 \chi_1(\Delta_{11})&\chi_1(\Delta_{59})\\
 \chi_2(\Delta_{11})&\chi_2(\Delta_{59})
 \end{pmatrix}
                  =\begin{pmatrix}0&2\\2&1\end{pmatrix}.
\tag{A.184.9}
\]
Its determinant is nonzero modulo three.  Pairing the
principal relation A.184.6, using A.184.8, gives this
matrix times \((k_{11,1},k_{59,1})^t\), equal to zero.
Both coefficients vanish.  A.184.7 now gives residue one
at both pure quadratic branches, and in particular
\(\theta_{n+1}(u)=0\).  This proves the second assertion
of A.184.2 for every actual relative norm-one global unit.
The proof has used its zero divisor at every prime; no
statement about a larger relative \(S\)-unit kernel is
inferred from A.184.6.

\emph{The secondary map and two consecutive norm equations.}
Every \(u\in M_n\) has an actual global-unit preimage by
the first part.  Two preimages differ by a relative
norm-one global unit.  The result just proved makes their
\(\theta_{n+1}\) values identical.  Multiplication of
preimages proves that \(\psi_n\) is a homomorphism.
A unit \(u\in M_n\) is a norm from two steps precisely
when its first preimage can be chosen in
\(\ker\theta_{n+1}\).  The value is independent of the
choice, and the first equality of A.184.2 at level \(n+1\)
identifies that kernel with its one-step global-unit norm
image.  This proves the fourth line of A.184.2 in both
directions.

For a scalar \(u\in\mathcal U_n\), its image in
\(\mathcal U_{n+1}\) has norm \(u^3\), and its two
residues are unchanged because both changed-prime residue
degrees are one.  Hence the defining formula gives
\(\psi_n(u^3)=\theta_{n+1}(u)=\theta_n(u)\), with the
fixed generators in A.184.1.

\emph{The completed domains and torsion maps.}
The pure fields have odd degree over \(\Q\) and are
subfields of \(\R\), so their roots of unity are
\(\{1,-1\}\), and norm on this subgroup is the identity.
Writing
\(V_n=\ker(N_{K_n/\Q}:\mathcal U_n\to\{1,-1\})\)
gives
\[
 \mathcal U_n=\{1,-1\}\times V_n,
 \qquad V_n\text{ free of finite rank},\qquad
 \widehat{\mathcal U}_{n,3}=V_n\otimes\Z_3.
\tag{A.184.10}
\]
Both tame characters kill \(-1\), and \(\psi_n\) does
also, by its preimage \(-1\).  The norm maps on the free
lattices are integral linear maps.  Flatness of \(\Z_3\)
preserves their kernels and images.  In particular their
completed norm-one kernel is the completion of the actual
kernel, and the vanished character remains zero there.
The finite target \(\F_3^2\) and the finite-index inclusion
\(M_n\subset\mathcal U_n\) give
\[
 \widehat M_{n,3}=\ker\widehat\theta_n,
 \qquad
 N_{n+1,n}\widehat{\mathcal U}_{n+1,3}
                         =\widehat M_{n,3}.
\tag{A.184.11}
\]
The secondary map extends uniquely to this compact group.
The same exact-lattice calculation, or the preimage proof
above applied to these compact groups, gives every completed
equality in A.184.2.  This explicitly retains the completed
source, rather than replacing it by arbitrary local classes.

\emph{A fixed scalar obstruction at every height.}
The exact residues of A.181.6 give
\(\theta_n(U)=(2,2)\) at every \(n\).
The scalar norm proves \(U^3\in M_n\), while
\[
                  \psi_n(U^3)=(2,2)\ne0.
\tag{A.184.12}
\]
Thus \(U^3\) is not a two-step global-unit norm, including
after completion.  Its cube is the literal scalar norm
\(N_{K_{n+2}/K_n}(U)=U^9\), so the finite unit support
class has exact order three in the actual kernel
\[
 \ker\left(\mathcal U_n/N_{n+2,n}\mathcal U_{n+2}
       \longrightarrow K_n^\times/N_{n+2,n}K_{n+2}^\times\right)
\tag{A.184.13}
\]
and in the kernel of its completed comparison map.  Membership
in these kernels follows from the field-norm verification next.

For completeness the field-norm and inverse-fibre maps from
A.181 apply to this shifted base with no change of radicand
transition.  At every \(m>n\) they use the actual groups
\[
 \operatorname{Gal}(\Omega_m/K_n)
      =(3^n\Z/3^m)\rtimes(\Z/3^m)^\times,
 \qquad\operatorname{Gal}(\Omega_m/K_m)
      =(\Z/3^m)^\times.
\tag{A.184.14}
\]
The local valuations of \(\rho_n\) are one at three
and eight at each changed-prime branch.  At the quadratic
branches \(U^3\) has cube residue, which is a
\(3^{m-n}\)-power residue for every \(m>n\) because
the residue three-part has order three.  Linear changed
branches impose no unit condition; at three the pure norm
is onto; elsewhere the unit root algebras are unramified.
The full decomposition subgroup in A.184.14 gives the
same exact global Hasse norm injection, and compact unit
fibres give compatible local families.
In a coherent global field family, all valuations at the
unique total-ramification branches would be zero.  At every
other \(K_n\)-place the actual affine Frobenius multiplier
is its residue cardinality.  The fixed-point and uniform
cycle-growth argument A.181.11--13 forces all its inverse
divisor coordinates to be zero as well.  The completed
principal-divisor injection in that argument remains the
finite-class-index injection at each of these same fields.
Thus a coherent actual or completed field family over
\(U^3\) would be a global-unit family.  Its two-step
projection contradicts A.184.12, proving both empty inverse
fibres at the shifted base.
\end{proof}

The executed certificate
\path{code/two_inert_secondary_unit_character_certificate.py}
retains the full finite ramified-residue calculation, the
anti-invariant class-character matrix, actual uniformizer
and scalar-norm exponents, and the exact norm/secondary-map
identities.  The all-height conclusion uses the displayed
ideal extension, norm, Hilbert 90 and fixed-divisor maps.
